% Auto-extracted during the research cleanup pass. Original source is preserved unchanged.
% Source: innerDense.tex, lines 2044--2535. Preserved exploration; not included in the main spine.

\subsection{Prospective block-recursive BCH inner}
\label{subsec:block-recursive-bch-inner}

This subsection records a possible replacement for the scalar fixed-tap dense inner.  It is not used in the proved
theorems above.  The point is to keep the recursive spreading mechanism, but replace the one-bit fixed-tap termination
rule by a short block-code expansion/collision rule.

Assume for simplicity that \(N=Mb\).  Partition the interleaved input into blocks
\[
U=(U_1,\ldots,U_M),\qquad U_i\in\F_2^b.
\]
Let \(C\subseteq\F_2^{2b}\) be a binary \([2b,b,d_0]\) code with a systematic generator written in
output-state coordinates as
\[
C=\{(v,Pv):v\in\F_2^b\},
\qquad
P:\F_2^b\to\F_2^b.
\]
For example, \(C\) could be an extended BCH or RM block code after choosing a systematic coordinate half.  The preferred
rate-one block recursion includes fixed state-coordinate permutations between blocks:
\begin{equation}
\label{eq:block-recursive-bch-inner}
V_i := U_i+S_{i-1},\qquad
Y_i := V_i,\qquad
S_i := \pi_i(P V_i),
\end{equation}
with \(S_0=0\), where the \(\pi_i\)'s are public fixed permutations of the \(b\) state coordinates.  Thus the current
mismatch \(V_i\) is emitted immediately, while a permuted copy of the parity half \(PV_i\) is stored as the next
feedback state.  Since \((V_i,PV_i)\in C\), every nonzero branch satisfies
\[
\wt(Y_i)+\wt(S_i)=\wt(V_i)+\wt(PV_i)\ge d_0.
\]
The orientation in \eqref{eq:block-recursive-bch-inner} is intentional.  A systematic encoder of the opposite form
\((Pv,v)\) gives dense immediate output but leaves a sparse feedback state when \(v\) is sparse, making a later
cancellation too cheap.  In the output-systematic orientation, a sparse input burst may emit only a few symbols in the
first block, but the stored state is the parity expansion \(Pv\), which is the object that the next input block must
hit exactly to turn the episode off.

\paragraph{State permutations and weight-only transitions.}
The state permutations in \eqref{eq:block-recursive-bch-inner} are not a rate cost, and for \(b=64\) they are only a
small fixed bit permutation per block.  They also remove low-power algebraic resonances of the deterministic map \(P\).
For the analysis we treat the \(\pi_i\)'s as independent uniform permutations, or equivalently prove that a typical
fixed draw has the desired finite certificate.  The update is
\begin{equation}
\label{eq:block-recursive-virtual-perm}
S_i := \pi_i(PV_i),
\end{equation}
where each \(\pi_i\) is uniform on the \(b\) coordinates.  After conditioning on \(\wt(S_i)\), this lets us track only
the state weight, rather than the exact support of \(S_i\), when upper bounding the probability that \(U_{i+1}=S_i\).
This is the same weight-only transition law that the earlier virtual-permutation analysis suggested, but now the
randomization is part of the candidate construction rather than only a proof convenience.

The local statistic needed for this analysis is the parity-expansion profile
\begin{equation}
\label{eq:block-recursive-expansion-profile}
B_{j,q}
:=
\#\{v\in\F_2^b:\wt(v)=j,\ \wt(Pv)=q\}.
\end{equation}
This replaces the ordinary local weight enumerator for the inner proof.  A branch with \(\wt(V_i)=j\) emits \(j\)
ones immediately and creates a state of weight \(q\) with multiplicity \(B_{j,q}\).

\begin{lemma}[One-block interleaver law against a fixed state]
\label{lem:block-recursive-one-block-law}
Expose any history that is independent of the labels in the next \(b\)-block, and suppose that \(h_{\rm rem}\) input
ones remain to be placed among \(N_{\rm rem}\) unexposed coordinates, including the next block.  Let
\(S\subseteq[b]\) be any fixed state support of size \(q\).  If \(U\) denotes the next input block, then for every
\(r,\ell\),
\begin{equation}
\label{eq:block-recursive-overlap-law}
\PP[\wt(U)=r,\ |U\cap S|=\ell]
=
\frac{\binom{q}{\ell}\binom{b-q}{r-\ell}\binom{N_{\rm rem}-b}{h_{\rm rem}-r}}
{\binom{N_{\rm rem}}{h_{\rm rem}}}.
\end{equation}
In particular,
\begin{equation}
\label{eq:block-recursive-collision-prob}
\PP[U_{i+1}=S_i]
\;=\;
\frac{\binom{N_{\rm rem}-b}{h_{\rm rem}-q}}{\binom{N_{\rm rem}}{h_{\rm rem}}},
\end{equation}
with the usual convention that impossible binomial coefficients vanish.  Thus turn-off is controlled by the random
interleaver collision with the current state, not by an additional dense random tap.
\end{lemma}

\begin{proof}
Choose the \(h_{\rm rem}\) remaining input positions uniformly from the \(N_{\rm rem}\) unexposed coordinates.  To have
\(\wt(U)=r\) and \(|U\cap S|=\ell\), choose \(\ell\) positions from \(S\), \(r-\ell\) positions from the other
\(b-q\) coordinates of the block, and the remaining \(h_{\rm rem}-r\) positions outside the block.  This gives
\eqref{eq:block-recursive-overlap-law}.  The exact-collision event is the special case \(r=\ell=q\).
\end{proof}

\begin{corollary}[Virtual permutations give a weight-only transition]
\label{cor:block-recursive-virtual-weight-law}
In the virtual-permuted recursion \eqref{eq:block-recursive-virtual-perm}, conditional on the current state weight
\(q\), the joint law of
\[
\wt(U_i),\qquad |U_i\cap S_{i-1}|,\qquad \wt(V_i)=\wt(U_i+S_{i-1})
\]
depends only on \(q\), \(b\), \(h_{\rm rem}\), and \(N_{\rm rem}\), not on the support labels of \(S_{i-1}\).  Moreover,
conditional on \(\wt(V_i)=j\), the support of \(V_i\) is uniform over the \(j\)-subsets of the block.
\end{corollary}

\begin{proof}
The first statement is immediate from Lemma~\ref{lem:block-recursive-one-block-law}, since
\(\wt(U_i+S_{i-1})=r+q-2\ell\).  For the uniformity statement, fix any two \(j\)-subsets \(A,A'\subseteq[b]\).  A block
coordinate permutation maps \(A\) to \(A'\) and preserves the law of the virtual state and of the next input block under
the uniform interleaver exposure.  Hence the conditional probabilities of \(V_i=A\) and \(V_i=A'\), given
\(\wt(V_i)=j\), are equal.
\end{proof}

\begin{corollary}[Branch expansion through the local profile]
\label{cor:block-recursive-profile-transition}
Under the virtual-permuted analysis model, conditional on \(\wt(V_i)=j\),
\[
\PP[\wt(S_i)=q\mid \wt(V_i)=j]
=
\frac{B_{j,q}}{\binom{b}{j}}.
\]
Consequently a weight-only episode ledger may use the transition
\[
q_{\mathrm{old}}
\longrightarrow
(r,\ell,j=q_{\mathrm{old}}+r-2\ell)
\longrightarrow
q_{\mathrm{new}},
\]
where the first arrow is governed by \eqref{eq:block-recursive-overlap-law} and the second by
\eqref{eq:block-recursive-expansion-profile}.
\end{corollary}

\begin{proof}
By Corollary~\ref{cor:block-recursive-virtual-weight-law}, conditional on \(\wt(V_i)=j\), the mismatch support is
uniform among all \(\binom{b}{j}\) \(j\)-subsets.  The number of such supports whose parity image has weight \(q\) is
exactly \(B_{j,q}\).
\end{proof}

\begin{corollary}[Heavy states cannot be canceled by too little remaining input]
\label{cor:block-recursive-heavy-state-charge}
At the start of any block, suppose the current state has weight \(q\) and at most \(H\) input ones remain in the whole
unexposed suffix.  Then the next emitted block satisfies
\[
\wt(Y_i)=\wt(U_i+S_{i-1})\ge (q-H)_+.
\]
In particular, if \(q>H\), exact turn-off in the next block is impossible.
\end{corollary}

\begin{proof}
Let \(r=\wt(U_i)\).  Since \(r\le H\), the symmetric difference of \(U_i\) and \(S_{i-1}\) has size at least
\(q-r\ge q-H\).  Exact turn-off would require \(U_i=S_{i-1}\), hence \(r=q\), impossible when \(q>H\).
\end{proof}

\paragraph{Dynamic-programming certificate form.}
The previous corollaries give an exact finite-state law for the virtual weight process before imposing the distance
budget.  If \(H\) input ones remain among \(R\) unexposed coordinates at the start of a block and the current state has
weight \(q\), define
\[
K_{R,H}^{(b,P)}(q;r,\ell,j,q')
:=
\frac{\binom{q}{\ell}\binom{b-q}{r-\ell}\binom{R-b}{H-r}}{\binom{R}{H}}\,
\mathbf{1}_{j=q+r-2\ell}\,
\frac{B_{j,q'}}{\binom{b}{j}}.
\]
This kernel simultaneously records the number \(r\) of input ones consumed in the block, the overlap \(\ell\) with the
current state, the emitted weight \(j=\wt(Y_i)\), and the next state weight \(q'=\wt(S_i)\).  Thus a certificate for
\[
p_h^{\mathrm{block\mbox{-}rec}}(\delta)
=
\PP[\wt(Y_1,\ldots,Y_M)\le d\mid \wt(U)=h]
\]
can be organized by the recursion
\[
F_i(H,q,D)
\le
\sum_{r,\ell,j,q'}
K_{R_i,H}^{(b,P)}(q;r,\ell,j,q')\,
F_{i+1}(H-r,q',D-j),
\]
with \(R_i=(M-i+1)b\), and impossible or negative-budget states set to zero.  The terminal condition depends on the
chosen boundary convention.  If a terminal state cost \(T(q)\) is charged, take
\[
F_{M+1}(0,q,D)=\mathbf{1}_{D\ge T(q)}.
\]
Tail-biting or forced termination corresponds to \(T(0)=0\) and \(T(q)=+\infty\) for \(q>0\); an explicit flush tail is
modeled by appending the corresponding zero-input transfer blocks.  This is the finite analogue of the scalar episode
ledger, but the termination/collision probabilities are now the hypergeometric terms in \(K\), while the local
spreading is captured by \(B_{j,q'}\).

There is one boundary convention that must be fixed before this becomes a theorem for a finite code.  The final state
\(S_M\) is not emitted by \eqref{eq:block-recursive-bch-inner}.  One can handle this by appending a short deterministic
flush/tail, by using a tail-biting block recursion, or by adding \(S_M\) to the distance ledger as an explicit terminal
state cost.  The analysis should not simply ignore \(S_M\), since doing so would create the same kind of artificial
late-tail loophole that appeared in the scalar finite-\(N\) diagnostics.

\paragraph{Proof target.}
The desired small-weight inner theorem would combine \eqref{eq:block-recursive-expansion-profile} and
\eqref{eq:block-recursive-collision-prob} in an episode ledger.  Each active block either:
\begin{itemize}
\item emits the current mismatch weight \(j\) and stores a state of weight \(q\);
\item terminates in the next block only if the interleaver places exactly that \(q\)-set as \(U_{i+1}\);
\item or survives, in which case the next mismatch is \(U_{i+1}+S_i\) and the process repeats.
\end{itemize}
The finite-\(N\) bottleneck should then be governed by expansion of low-weight words under \(P\) and by the exact
hypergeometric collision factors above.  For an extended BCH \([128,64,22]\)-shaped block, initial diagnostics of the
output-systematic parity half give \(\wt(Pv)\) near \(32\) on average for \(\wt(v)=1,2,3,4\), with sampled minima in
the high teens or above.  This is only evidence, but it suggests that the recursion may suppress the tiny-weight
turn-off modes that dominate the scalar fixed-tap finite-\(N\) certificate.  Indeed, once a tiny mismatch creates a
state with \(\wt(Pv)>h_{\rm rem}\), Corollary~\ref{cor:block-recursive-heavy-state-charge} says that the episode cannot
turn off immediately; the excess state weight must be emitted in later blocks or charged at the boundary.

The first finite probe supports this interpretation.  With \(N=2^{21}\), \(b=64\), the output-systematic extended BCH
\([128,64,22]\) parity half, and terminal-state weight charged at the end, the exact \(h=1\) failure probability is just
the block-level late-start effect:
\[
\begin{array}{c|c|c|c}
\delta & d=\lfloor\delta N\rfloor & p_1^{\mathrm{block\mbox{-}rec}}(\delta) & \log_2 p_1\\
\hline
.106 & 222298 & .212075233 & -2.237352\\
.09 & 188743 & .180082321 & -2.473272
\end{array}
\]
In contrast, a Monte Carlo probe at \(\delta=.106\) saw no failures in \(3000\) trials at \(h=22\), and no failures in
\(1000\) trials for \(h\in\{8,16,22,24,25,32,40,64\}\).  This is not a certificate, but it is qualitatively different
from the scalar fixed-tap bottleneck: the dangerous mass appears to collapse to genuinely tiny \(h\), while the
old \(h\approx22\text{--}32\) ridge is not visible in the probe.  The diagnostic entry point is
\[
\texttt{python scripts/probe\_block\_recursive\_inner.py}.
\]

The same probe also warns against overreading the zero-failure samples.  Once turn-off is suppressed, the remaining
first-order obstruction is ordinary block-level late placement.  If a live zero-input trajectory emits about \(b/2=32\)
ones per block, then a support of weight \(h\) can still fit below \(d\) when its first active block lies in roughly the
last \(\lfloor d/32\rfloor\) block positions.  The crude placement factor
\[
\frac{\binom{64\lfloor d/32\rfloor}{h}}{\binom{N}{h}}
\]
therefore gives a useful calibration.  At \(N=2^{21}\) and \(\delta=.106\), this factor contributes about
\(-71.62\) bits at \(h=32\).  Paired with the RM block-outer \(h=32\) coefficient (\(\log_2 A_{32}^{\rm out}\approx
37.63\)), the resulting single-term margin is only about \(34\) bits.  At \(\delta=.09\), the same calculation gives
about \(41.54\) bits.  Thus the block-recursive inner appears to remove the scalar termination ridge, but a \(40\)-bit
\(\delta=.106\) finite checkpoint would still need a slightly stronger outer floor/spectrum, a sharper output-growth
certificate, or a smaller target distance.  The late-start balance diagnostic is
\[
\texttt{python scripts/estimate\_block\_recursive\_late\_start.py}.
\]

The sharper version of this calibration is the zero-input column profile
\[
D_L(P):=\min_{s\ne0}\sum_{t=0}^{L-1}\wt(P^t s).
\]
For the small exact \([32,16,8]\) toy, \(D_L(P)\) initially follows the two-block bound \(D_L\ge\lfloor L/2\rfloor d_0\)
but improves for longer windows.  For the \([128,64,22]\) candidate, exact enumeration over all weight-one and
weight-two initial states gives the largest live zero-input windows that still fit under the target budgets:
\[
\begin{array}{c|c|c|c}
\wt(s_0) & \delta & \max\{L:D_L(s_0)\le \lfloor\delta N\rfloor\} & \log_2(L/32768)\\
\hline
1 & .106 & 6978 & -2.231402\\
2 & .106 & 7003 & -2.226243\\
1 & .09 & 5930 & -2.466184\\
2 & .09 & 5949 & -2.461569
\end{array}
\]
Thus the actual low-weight column growth is close to the crude \(32\)-per-block late-start estimate, not to the weak
pair bound.  This supports a concrete proof target: certify a usable column-distance/profile lower bound for the
reachable low-weight state classes, then add the hypergeometric cancellation terms from
Lemma~\ref{lem:block-recursive-one-block-law}.  The column-profile diagnostics are
\[
\texttt{python scripts/profile\_block\_recursive\_column.py}
\qquad\text{and}\qquad
\texttt{python scripts/profile\_block\_recursive\_late\_threshold.py}.
\]

Combining the exact low-weight threshold with the exact first-active-block occupancy law isolates the first realistic
finite certificate target.  If a weight-\(h\) input is uniform over \(M=N/b\) blocks, the probability that its first
active block lies in the last \(L\) blocks is
\[
\sum_{r=1}^h\sum_{t=M-L+1}^{M}
\frac{\binom{b}{r}\binom{b(M-t)}{h-r}}{\binom{N}{h}},
\]
where \(r\) is the occupancy of the first active block.  For the RM block outer, the first nonzero outer weight is
\(h=32\).  At \(\delta=.09\), the exact weight-one/two column thresholds above give \(L=5949\) blocks, and the
first-active late probability has logarithm \(-78.771731\).  Multiplying by the RM \(h=32\) outer coefficient
\(\log_2 A_{32}^{\rm out}=37.629712\) leaves \(41.142\) bits of margin before handling the remaining branches.
At \(\delta=.106\), the analogous \(L=7003\) calculation leaves only \(33.611\) bits at \(h=32\).  Thus the clean
first target is a finite \(\delta=.09\) block-recursive certificate: prove that the early-first-active branch cannot
recover enough distance through later hypergeometric cancellations, then sum the already-negative late branch against
the RM outer spectrum.

A conditioned early-branch probe shows that this split should include a small buffer.  At \(\delta=.09\), \(h=32\),
and \(L=5949\), forcing the first active block to be exactly one block before the late threshold still produced rare
Monte Carlo failures, because the zero-input trajectory is only barely above the distance budget and later input blocks
can cancel the first trajectory.  By ten blocks before the threshold no failures appeared in the same small probe, and
by \(100\) blocks before the threshold the zero-input excess was already several thousand output symbols.  Thus the
certificate should use a buffered late window
\[
L_{\mathrm{late}}=5949+B
\]
and prove the complement only after this buffer.  The RM \(h=32\) margin remains about \(40.75\) bits at
\(L_{\mathrm{late}}=6000\) and about \(40.36\) bits at \(L_{\mathrm{late}}=6050\), while \(L_{\mathrm{late}}=6100\) is
essentially break-even for a \(40\)-bit margin.  Since the \(40\)-bit target is only a calibration point, the important
lesson is structural: the finite proof should have three branches, not two--a late/buffer branch, an early branch with
substantial column-distance excess, and a hypergeometric cancellation ledger for the buffer boundary.  The conditioned
diagnostic is
\[
\texttt{python scripts/probe\_block\_recursive\_early\_branch.py}.
\]

For the actual first-moment sum one should not quote only the \(h=32\) peak term.  Summing the exact late/buffer
probability against the RM outer prefix through \(h=500\) gives
\[
\begin{array}{c|c|c}
L_{\mathrm{late}} & \log_2\sum_{h\le 500}A_h^{\rm RM}\Pr[\text{first active block in final }L_{\mathrm{late}}]
& \text{margin bits}\\
\hline
5949 & -40.115431 & 40.115431\\
6000 & -39.708713 & 39.708713\\
6050 & -39.313105 & 39.313105
\end{array}
\]
Thus the late/buffer branch is already in certificate form, but the full prefix is roughly one bit worse than the
single \(h=32\) term.  The remaining proof obligation is entirely in the complement: either show that the early branch
has enough zero-input column-distance surplus, or charge a boundary failure to a large hypergeometric overlap between
the later input trajectories and the first trajectory.  A deliberately pessimistic diagnostic for this last step is
\[
\texttt{python scripts/bound\_block\_recursive\_cancellation.py},
\]
which computes exact trajectory-overlap populations and applies an exponential-moment bound; its current role is to
size the cancellation lemma, not yet to serve as the final certificate.

There is also a useful bookkeeping simplification.  Instead of proving a delicate cancellation statement immediately at
the boundary, one may simply union-bound a short strip before the exact \(L=5949\) threshold.  For the same RM prefix,
the strip \(1\le g\le 51\) blocks before the threshold has total first-moment mass \(2^{-41.733963}\), and the strip
\(1\le g\le 101\) has total mass \(2^{-40.542230}\).  The peak contribution in both cases is still \(h=32\),
first-block occupancy \(r=1\), at the far edge of the strip.  Thus a clean finite split is:
\[
\begin{array}{c}
\text{exact late threshold}\\
\text{short boundary strip by union bound}\\
\text{early branch with column-distance surplus}
\end{array}
\]
The strip diagnostic is
\[
\texttt{python scripts/certify\_block\_recursive\_boundary\_strip.py}.
\]

The first cancellation ledger also changes the proof target.  A naive bound that treats an exact turnoff as immediate
failure is much too pessimistic, because the remaining input weight must still generate a low-weight continuation.  In
fact the dominant exact-turnoff row for \(h=32,r=1\) is a very short pair: after two zero-input BCH steps, some
weight-one states return to weight one, and a single later input can turn the episode off after only about \(36\)
output symbols.  This is not by itself a bad codeword, but it means the early-branch proof must count short
cancellation chains or clusters.  The low-state recurrence profile begins
\[
\begin{array}{c|c|c|c|c}
\wt(s_0) & t & \min \wt(P^t s_0) & \mathbb{E}\wt(P^t s_0) &
\#\{s_0:\wt(P^t s_0)\le 1\}\\
\hline
1&2&1&15.562500&30\\
1&4&1&23.875000&16\\
1&6&1&28.093750&7\\
1&8&1&30.906250&2
\end{array}
\]
and there are no weight-one returns among unit states for \(10\le t\le 12\) in the same diagnostic.  Thus the revised
early theorem should be an episode/cluster theorem: supports with many short cancellation edges are combinatorially
rare under the interleaver, while supports without such a cluster accumulate the column-distance surplus.  The relevant
diagnostics are
\[
\begin{array}{c}
\texttt{python scripts/certify\_block\_recursive\_exact\_turnoff.py}\\
\texttt{python scripts/profile\_block\_recursive\_low\_state\_recurrence.py}
\end{array}
\]

This short-return obstruction appears to be an artifact of iterating the same parity map \(P\).  With independent
state-coordinate permutations between BCH blocks, the trajectory is no longer \(P^t s_0\) but
\[
S_t=\pi_t P\pi_{t-1}P\cdots \pi_1P s_0.
\]
In a small diagnostic with \(20\) random permutation schedules and \(t\le 12\), every unit start became dense after the
first step and no unit trajectory returned to weight at most \(16\).  Weight-two starts also stayed essentially dense:
there were no returns to weight at most \(8\), and only isolated returns to weight at most \(16\) among the \(2016\)
weight-two starts.  This is the main reason to make the \(\pi_i\)'s part of the construction rather than merely a
virtual proof device.  A full Monte Carlo probe with the same state permutations still sees \(h=1\) failures from
ordinary late placement, but in a small \(\delta=.09\) run saw no \(h\ge 2\) failures except rates consistent with
late placement for \(h=2\).  The permutation recurrence diagnostic is
\[
\texttt{python scripts/profile\_block\_recursive\_permuted\_recurrence.py}.
\]

With the state permutations included, the natural proof object is a finite weight-chain certificate.  Conditional on
the current state weight \(q\), the current block occupancy \(r\), and the overlap \(\ell\), the emitted weight is
\[
j=q+r-2\ell,
\]
and the next state weight is distributed according to \(B_{j,q'}/\binom{b}{j}\).  Thus one can upper-bound a continuation
probability by the Chernoff transfer
\[
\Pr[\wt(Y_{\rm cont})\le d_{\rm rem}]
\le
\exp(\lambda d_{\rm rem})
\mathbb{E}\exp\{-\lambda\wt(Y_{\rm cont})\},
\]
where the expectation is over the exact without-replacement block occupancies and the local BCH profile \(B_{j,q'}\).
This is the theorem-shaped replacement for the deterministic \(P^t\) column-profile proof.  The present Python
prototype,
\[
\texttt{python scripts/certify\_permuted\_inner\_chernoff.py},
\]
is only a skeleton: exact profiling for all \(j\) is impossible by enumeration, and the current conservative fallback
for unprofiled \(j\)'s is too loose near the finite \(\delta=.09\) threshold.  The next certificate step is therefore
to replace the fallback by a rigorous local BCH lower-tail bound for \(B_{j,q'}\), or to implement the same transfer
with a coarser but certified local profile.

For proof development it is useful to separate this missing local theorem as a hypothesis.  Assume that for every
\(j\) there is an explicit envelope \(\Phi_b(j,q)\) such that
\[
\Pr_{v:\wt(v)=j}\big[\wt(Pv)\le q\big]\le \Phi_b(j,q),
\qquad
\Phi_b(j,q)=0\quad(q<d_0-j).
\]
The working guess is the BCH floor plus a random-vector lower tail with slack,
\[
\Phi_b(j,q)\approx
\min\left\{1,\;2^{s}\Pr[{\rm Bin}(b,1/2)\le q]\right\},
\]
where \(s\) is a slack budget to be determined.  Conditional on this local hypothesis, the rest of the inner proof is
finite and mechanical: convert \(\Phi_b\) into a pessimistic next-state distribution by placing as much mass as allowed
on small \(q'\), run the Chernoff transfer, and sum the resulting continuation bound against the outer spectrum and the
first-active-block placement law.

The first conditional runs with the zero-slack random-tail guess and coarse \(16\)-weight state buckets show the shape
of the remaining finite problem.  For \(h=32,r=1,q\approx32\) at \(N=2^{21}\), \(\delta=.09\), and remaining distance
about \(188710\), the continuation Chernoff bound is still trivial at \(6050\) live blocks, buys only a few bits around
\(8000\) live blocks, and buys about \(15.5\) bits by \(10000\) live blocks:
\[
\begin{array}{c|c}
\text{live blocks after first} & \text{best observed }\log_2\text{ continuation bound}\\
\hline
6050 & +5.59\\
8000 & -2.66\\
10000 & -15.54
\end{array}
\]
Thus the near-boundary region cannot be ignored entirely, but farther early starts should be controllable once the
local envelope is made rigorous and the transfer is optimized.  The exact placement mass of a boundary strip also has a
closed form:
\[
\Pr[\text{first active block in the }G\text{-block strip before }L]
=
\frac{\binom{b(L+G)}{h}-\binom{bL}{h}}{\binom{N}{h}}.
\]
Against the RM prefix through \(h=500\), \(L=5949\) and \(G=1051\) give total mass \(2^{-32.326936}\), while
\(G=2051\) gives \(2^{-25.823633}\).  These are useful bookkeeping numbers, but they are not enough on their own; the
conditional continuation certificate must be included for the longer early range.

The first bucketed conditional first-moment check is encouraging in the limited sense we need at this stage.  Using the
zero-slack random-tail envelope, \(16\)-weight state buckets, and the convention that each gap bucket uses the
continuation bound at its near edge, the dominant \(h=32,r=1\) branch over the whole non-late range has total exponent
roughly in the mid-teens below zero.  A coarse split gives
\[
\begin{array}{c|c|c}
\text{gap range} & \log_2\text{ placement+continuation+outer term} & \text{near-edge live blocks}\\
\hline
1\text{--}12000 & -16.17 & 5949,9949,13949\text{ buckets}\\
12001\text{--}19000 & -17.25 & 17949\\
19001\text{--}26819 & -19.39 & 24949
\end{array}
\]
For \(r=2\), the same \(1\text{--}12000\) coarse split gives \(-25.48\) bits, and for \(r=3,4\) it gives
\(-35.38\) bits combined.  Finally, \(r\ge5\) has total exponent \(-9.79\) by placement alone across the whole
non-late range.  Therefore the guessed local envelope appears strong enough to make the \(h=32\) ridge negative, even
with very coarse bucketing.  This is not yet a certificate: the final version must refine or justify the bucket
monotonicity convention, include all relevant \(h\), and replace the guessed \(\Phi_b\) by a proved local BCH bound.
The bucketed diagnostic is
\[
\texttt{python scripts/certify\_permuted\_inner\_bucket\_sum.py}.
\]

The first nearby-\(h\) probes suggest that the \(h\)-summation still deserves real attention.  Under the same guessed
zero-slack random-tail envelope and \(16\)-weight state buckets, \(h=40,r=1\) remains negative on the three coarse
gap buckets \(1\text{--}4000\), \(4001\text{--}8000\), and \(8001\text{--}12000\), with exponents approximately
\(-19.44\), \(-16.49\), and \(-24.11\).  However, the same middle bucket at \(h=64,r=1\) gives only about
\(-11.22\) bits, and at \(h=80,r=1\) it gives about \(-7.62\) bits.  Thus the conditional certificate is not merely a
tiny-\(h\) check: larger intermediate \(h\)'s may be closer to the ridge once the outer spectrum growth is included.
The current pure-Python transfer becomes slow in this regime.  Replacing huge binomial evaluations by a recurrence for
the block-occupancy hypergeometric law gives a useful speedup, but broad \(h\)-sweeps are still too slow.  The next
finite-certificate step is either to accelerate the weight-chain kernel, compress the rare high-occupancy block tail
with a certified upper bound, or add an analytic tail envelope in \(h\) before sweeping broadly.
